\documentclass[11pt]{article}
\usepackage{amsmath,amssymb,geometry}
\geometry{margin=22mm}
\begin{document}

\section*{P360: complete nonlinear metro regulation repair}
Stations $s=0,\ldots,9$ enumerate the ten directional service locations
of five physical stations. Trains $j=0,\ldots,J-1$ share the same operating line. The original timetable uses $J=6$ and the large tier extends the timetable to the input train count.
Each directional service location has the finite storage state used below.
All time is seconds, mass tonnes, energy kWh and auxiliary power kW.
Input contains full arrays $r^{pla},l^{pla},d^{pla},w^{run},w^{dwe},
\lambda,\beta$, first arrivals $a_{0j}$ and preceding departures $d_{s,-1}$.
Capacity is $C=1200$, storage capacity $E^{max}=500$ and initial energy
$E^0=250$ at each modeled service location; these are explicit completed
scenario values supplied by the instance.

Choose controls $x_{sj}\in[-30,30]$ for $s<9$ and
$g_{sj}\in[-15,15]$. Define
\begin{align}
r_{sj}&=r^{pla}_{sj}+x_{sj}+w^{run}_{sj},&
l_{sj}&=l^{pla}_{sj}+g_{sj}+w^{dwe}_{sj},\\
d_{sj}&=a_{sj}+l_{sj},&
a_{sj}&=d_{s-1,j}+r_{s-1,j}\quad(s>0),\\
h_{sj}&=d_{sj}-d_{s,j-1}.&&
\end{align}
Require $h_{sj}\ge270$ for $j>0$, $h_{s0}\ge0$, $l_{sj}\ge20$,
$300\le r_{4j}\le600$, and the passenger dwell condition below.
The turnaround interval is the zero-distance transition between the two
directional service sequences. Controls may be negative, correcting the
original extracted blanket nonnegative domain.

\paragraph{Complete passenger dynamics.}
Put $n_{-1,j}=0$ and $m_{s,-1}=0$. Then
\[
 q_{sj}=\beta_{sj}n_{s-1,j},\quad
 p_{sj}=\lambda_{sj}h_{sj}+m_{s,j-1},\quad
 b_{sj}=\min\{p_{sj},C-n_{s-1,j}+q_{sj}\},
\]
\[
 n_{sj}=n_{s-1,j}+b_{sj}-q_{sj},\quad
 m_{sj}=p_{sj}-b_{sj},\quad
 l_{sj}\ge3+0.03(b_{sj}+q_{sj}).
\]
All counts here are continuous passenger-flow approximations;
$0\le n\le C$ and $b,q,p,m\ge0$. Left-behind passengers persist in the queue.
Passengers still in the network at the horizon remain in the terminal state.

\paragraph{Running level and energy tables.}
For each segment, input table breakpoints $\ell_{s\eta}$ increase strictly;
choose $\eta(s,j)=\max\{\eta:r_{sj}\ge\ell_{s\eta}\}$.
The last interval extends to the largest feasible running time;
the input lower breakpoint covers all feasible controls.
Let $M_{sj}=200+0.075n_{sj}$.
With the table's traction/braking coefficients $c^{tra}_{s\eta},
c^{brk}_{s\eta}$ and acceleration/deceleration durations
$A_{s\eta},B_{s\eta}$, define for $s<9$
\[
 E^{tra}_{sj}=M_{sj}c^{tra}_{s,\eta(s,j)},\quad
 E^{brk}_{s+1,j}=M_{sj}c^{brk}_{s,\eta(s,j)},\quad
 t^{acc}_{sj}=A_{s,\eta(s,j)},\quad
 t^{dec}_{s+1,j}=B_{s,\eta(s,j)}.
\]
Set first-station braking/deceleration and last-station
traction/acceleration to zero.

\paragraph{Immediate regenerative flow.}
The opposite service location is $\sigma(s)=9-s$.
For each braking train $j$ and traction train $k$ define the actual overlap
\[
 O_{sjk}=\max\{0,\min(a_{sj},d_{\sigma(s),k}+t^{acc}_{\sigma(s),k})
 -\max(a_{sj}-t^{dec}_{sj},d_{\sigma(s),k})\}.
\]
Nonnegative delivered flow $z_{sjk}$ satisfies
\[
 z_{sjk}\le \rho_2 O_{sjk},\qquad
 U_{sj}:=\sum_kz_{sjk}\le\rho_1E^{brk}_{sj},\qquad
 V_{tk}:=\sum_{j}z_{\sigma(t),j,k}\le E^{tra}_{tk}.
\]
Input values are $\rho_1=0.7$ and $\rho_2=500/3600$ kWh/s.
Immediate-flow imports equal exports. The code chooses one feasible allocation
in lexicographic $(s,j,k)$ order, each time taking the minimum residual
braking supply, overlap bound and traction demand. This deterministic choice
is an energy-allocation heuristic.

\paragraph{Storage and auxiliary energy.}
Let $\eta_c=\eta_d=0.9$.
At every modeled service location, sequentially for all trains,
\begin{align}
E^{st1}_{s0}&=E^0,&
E^{st1}_{sj}&=E^{after}_{s,j-1}\quad(j>0),\\
E^{sav}_{sj}&=\min\{\eta_c\max(0,E^{brk}_{sj}-U_{sj}/\rho_1),
                      E^{max}-E^{st1}_{sj}\},\\
E^{st2}_{sj}&=E^{st1}_{sj}+E^{sav}_{sj},\\
E^{tak}_{sj}&=\min\{\eta_d E^{st2}_{sj},\max(0,E^{tra}_{sj}-V_{sj})\},\\
E^{after}_{sj}&=E^{st2}_{sj}-E^{tak}_{sj}/\eta_d.
\end{align}
All energy amounts are nonnegative; storage stays within $[0,E^{max}]$.
The saved-energy formula removes braking energy already supplied immediately,
and delivered discharge is divided by efficiency in the state balance.
For $s>0$,
$E^{aux}_{sj}=(50+0.11n_{s-1,j})(a_{sj}-a_{s-1,j})/3600$;
$E^{aux}_{0j}=0$.

\paragraph{Full objective.}
With input weights $(\xi_1,\ldots,\xi_6)=(0.8,1,1,1.5,10,0.5)$ minimize
\begin{align}
F_1&=\xi_1\sum_{sj}(d_{sj}-d^{pla}_{sj})^2
       +\xi_2\sum_{s,j>0}(h_{sj}-300)^2,\\
F_2&=\sum_{sj}\{\tfrac12\xi_3\lambda_{sj}h_{sj}^2
                         +\xi_4m_{s,j-1}h_{sj}\},\\
F_3&=\xi_5\sum_{sj}(E^{tra}_{sj}+E^{aux}_{sj}-V_{sj}-E^{tak}_{sj}),\\
F_4&=\xi_6\{\sum_{s<9,j}x_{sj}^2+\sum_{sj}g_{sj}^2\},\qquad
F=F_1+F_2+F_3+F_4.
\end{align}
Source pp.6--11,18 supplies train/passenger/energy/control structure.
The original02 omitted the four objectives and level/overlap definitions,
misstated boarding capacity and omitted storage initialization/discharge loss.
The above is the explicit separate completed physical interpretation,
using the completed passenger-flow definitions above.

The unchanged solver uses bounded coordinate search and the feasible
immediate-flow allocation described above. All ten service locations,
six trains and finite capacities are retained; the output is a feasible MINLP solution. Independent complete replay of the saved output verifies all
listed equations/domains and the four objective components.

\end{document}
